\chapter{Local Realization}\label{ch:local}

The variation considered so far happens inside one room: several viewers, several realizations, one screen. There is a second kind of variation that selective cinema makes newly possible, between rooms rather than within them. A work distributed as a scene model rather than a fixed sequence can be realized differently in different places, fitted to the theater, the language, the culture, even the performers. This chapter examines local realization and the ways it combines with selection inside the room.

\section{Two axes of variation}

Distinguish \emph{intertheater} variation, in which different venues show different realizations, from \emph{intratheater} variation, in which one venue shows several realizations to different viewers. Ordinary distribution already has the first in a crude form: dubbed and subtitled prints, regional cuts, censored versions. Selective cinema adds the second. The two can be combined, so that each venue offers its own local family, and a work becomes a family of families.

\section{Fitting the room}

The mildest form of local realization adapts the work to the physical theater. Screen size and shape, seating distance, and the height of the rows affect how an image reads. A close-up that dominates a small screen overwhelms a very large one; fine detail visible from the front rows is lost at the back. A work distributed as a scene model with declared camera rules can be rendered for each room: framing adjusted to the aspect ratio, camera distance to the typical viewing distance, pacing of dense visual information to the size of the screen.

This is variation that changes nothing a viewer is told. It adjusts image distance between theaters and leaves fact distance at zero, and it is uncontroversial except for one thing: it means that no two venues show quite the same frames. The work that critics review, archives preserve, and courts cite is no longer a single sequence. \Cref{ch:which-film} takes up what the work then is.

\section{Language}

The most common local variation is language, and it collides immediately with \Cref{ch:composite-sound}. Language is carried by sound. A multilingual room watching one picture needs either individual sound, one language per viewer, or a shared track in one language with the rest subtitled. Subtitles are a visual variation, and a selective screening can carry them as separate streams: the same picture with different subtitles for each language, selected by the glasses. This is cheap in image distance, because only the text differs, and adaptive gating keeps the glasses open whenever no subtitle is on screen.

Dubbing is harder. A dubbed line does not match the lips that were filmed speaking the original. Under shared picture and individual sound, the mismatch is the ordinary one of dubbed film. Under individual pictures, it could be removed: each language stream could show lips synchronized to its own language, re-rendered or re-performed. This is now technically possible with synthetic re-rendering of faces, and it raises the question of performer consent, taken up in \Cref{ch:versions}: an actor who spoke one line has not thereby agreed to be shown speaking another.

\section{Local performers}

A stronger form of local realization replaces the performers. A story's roles, blocking, and timing are held fixed, and each region casts its own actors, who perform the same scenes in their own language, with their own gestures, in settings adjusted to local architecture and custom. The result is more than dubbing, because the visible performers are local, and less than a remake, because the structure is shared and the local realizations remain synchronized with a common production.

This raises a question about performance. If several actors play one character in one structural role, the character becomes partly separable from any single embodiment. That separation is familiar from theater, where many actors play Hamlet, and unfamiliar in cinema, where a role is usually fixed to the face that played it. Local realization moves cinema toward theater in this respect. It also cuts the other way: a local performer may change the character substantially, and to call the result a ``substitution'' can understate what the performer contributed. Credit, in such a system, has to follow the realization, not only the structure.

\section{Translation, localization, rewriting}

Local variation comes in degrees, and the degrees differ in what they claim.

\emph{Translation} renders the same meaning in another language. It aims at zero fact distance.

\emph{Localization} replaces references that would not be understood: a joke about one country's politics with a joke about another's, a sign in one script with a sign in another, a food, a currency, a holiday. It aims at the same effect by different means, and it may change what is established, in small ways, to keep the effect.

\emph{Substitution} replaces performers and settings while keeping the structure.

\emph{Rewriting} changes the structure itself to fit the local culture: different relationships, different motives, different outcomes.

The further along this sequence a local realization goes, the less plausible it is to call one version the original and the others adaptations. A rewritten local realization is a different work that shares a skeleton with the others. Selective cinema does not create this problem; it has always existed in the relation between a film and its remakes. It makes the problem simultaneous, because the realizations are designed and released together, and it makes it visible, because a single venue may offer several.

\section{The template problem}

Local realization carries a risk that needs to be named. A work designed as a common structure with local surfaces presupposes that the structure is neutral: that its story, its relationships, its sense of what is funny or moving or just, can be carried into any culture by changing the faces and the signs. That presupposition is rarely true. A structure is written somewhere, by someone, and it carries the assumptions of its origin. Local realizations built on it import those assumptions while dressing them in local detail, and the local detail can make the import harder to notice.

The danger is a kind of cultural template: one culture's story with everyone else's faces as skins. It is not an argument against local realization. It is an argument that local realization should reach, where needed, as far as rewriting, and that the people who realize a work locally should have authority over its structure as well as its surface. A system that lets local creators change the faces but not the story has decided in advance whose story it is.

\section{Local families}

The two axes combine. A venue can offer, in one room, a local family: realizations in two languages, or a comic and dramatic pair performed by local actors, or a beginner and expert realization adapted to a local curriculum. Everything in Parts~I and~II applies within each venue, and local realization applies between them.

The combination also sharpens a point that runs through the book. Selective cinema makes it possible for people in different places, and people in the same room, to be shown different things under the same title. When the differences are declared, chosen, and made by people with standing to make them, this is an extraordinary enlargement of what a film can do. When they are not, it is a means of telling different audiences different things while each believes it saw the work. Part~VIII is about the difference.
